

The role of identity in politics is crucial to navigating the modern political landscape, dominated by social media and the internet. Classic political speech, often not accessible unless viewed live in person or on television is commonly accessible through the internet, thus providing greater access to political speech. A fairly recent development in identity politics over the 2010s and 2020s has been ontological security and the role it plays in shaping states' actions \Parencites{mitzenOntologicalSecurityWorld2006}{kocanIdentityOntologicalSecurity2023}{davidPolicingMemoryBosnia2019}. Although there has been an increase in analysis of ontological security, it has primarly focused on states such as the United States, which are primarly seperated by parties in regards to political statements \parencite{mitzenOntologicalSecurityWorld2006}. The aim of this project is to expand the discourse of ontological security to include states that are highly ethnically charged in order to further grasp the role that ontological security plays in state identity.

The Republic of Bosnia and Herzegovina (BiH) is one of few countries that has a political system that has safeguards to provide equal representation of the different ethnic groups throughout the country, but this provides additional complications to political discourse and national identity. The goal of this project is to determine whether political discourse in BiH maintains a level of cohesion acroos parties and ethnicities, thus following the assumption that states seek ontological security, or whether disagreement between the parties and ethnicities actively instills ontological insecurity. As well as focusing primarly on discourse by politicians, analysis convey the actions carried out as a result of this discourse and any roll that BiH's neighbors may play in BiH's ontological (in)security.



